\SourceLeaf{78}\begin{center}
{\Large VIERTER ABSCHNITT}\par\bigskip
{\small ÜBER}\par\bigskip
{\Large KONGRUENZEN ZWEITEN GRADES.}
\end{center}

\begin{center}\rule{.09\linewidth}{.25pt}\end{center}

\Needspace{10\baselineskip}
\begin{center}\small\emph{Quadratische Reste und Nichtreste.}\end{center}

\Article{94}

\textsc{Satz.} \emph{Nimmt man eine beliebige Zahl $m$ als Modul, so können von den Zahlen $0,1,2,3....m-1$ höchstens $\tfrac12m+1$, wenn $m$ gerade ist, und höchstens $\tfrac12m+\tfrac12$, wenn $m$ ungerade ist, einem Quadrat kongruent sein.}

\emph{Beweis.} Da die Quadrate kongruenter Zahlen kongruent sind, ist jede Zahl, die einem Quadrat kongruent sein kann, auch einem Quadrat kongruent, dessen Wurzel $<m$ ist. Es genügt also, die kleinsten Reste der Quadrate $0,1,4,9....(m-1)^2$ zu betrachten. Man sieht leicht, dass $(m-1)^2\equiv1$, $(m-2)^2\equiv2^2$, $(m-3)^2\equiv3^2$ usw. gelten. Wenn $m$ gerade ist, haben daher die Quadrate $(\tfrac12m-1)^2$ und $(\tfrac12m+1)^2$, $(\tfrac12m-2)^2$ und $(\tfrac12m+2)^2$ usw. dieselben kleinsten Reste; wenn $m$ ungerade ist, sind die Quadrate $(\tfrac12m-\tfrac12)^2$ und $(\tfrac12m+\tfrac12)^2$, $(\tfrac12m-\tfrac32)^2$ und $(\tfrac12m+\tfrac32)^2$ usw. kongruent. Somit können nur Zahlen, die einem der Quadrate $0,1,4,9....(\tfrac12m)^2$ kongruent sind, einem Quadrat kongruent sein, wenn $m$ gerade ist. Ist der Modul ungerade, so muss jede Zahl, die einem Quadrat kongruent ist, einem der Quadrate $0,1,4,9....(\tfrac12m-\tfrac12)^2$ kongruent sein. Es gibt also im ersten Fall höchstens $\tfrac12m+1$, im zweiten höchstens $\tfrac12m+\tfrac12$ verschiedene kleinste Reste. \mbox{\emph{Q. E. D.}}

\emph{Beispiel.} Nach dem Modul $13$ findet man für die Quadrate der Zahlen $0,1,2,3...6$ die kleinsten Reste $0,1,4,9,3,12,10$; danach kehren dieselben \SourceLeaf{79}in umgekehrter Reihenfolge wieder, $10,12,3$ usw. Jede Zahl, die keinem dieser Reste kongruent ist, also eine Zahl, die einer der Zahlen $2,5,6,7.\ 8,11$ kongruent ist, kann daher keinem Quadrat kongruent sein.

Nach dem Modul $15$ findet man die Reste $0,1,4,9,1,10,6,4$, nach denen dieselben in umgekehrter Reihenfolge wiederkehren. Hier ist die Anzahl der Reste, die einem Quadrat kongruent sein können, sogar kleiner als $\tfrac12m+\tfrac12$, denn es sind $0,1,4,6,9,10$. Die Zahlen $2,3,5,7,8,11,12,13,14$ und die Zahlen, die einer von ihnen kongruent sind, können keinem Quadrat nach dem Modul $15$ kongruent sein.

\Article{95}
Daraus folgt, dass man für jeden Modul alle Zahlen in zwei Klassen einteilen kann: Die eine enthält die Zahlen, die einem Quadrat kongruent sein können, die andere diejenigen, bei denen dies nicht möglich ist. Die ersteren nennen wir \emph{quadratische Reste der als Modul gewählten Zahl}\HistoricalNote{*)}{Genau genommen verwenden wir hier den Genitiv in einem anderen Sinn als bisher. Wir müssten sagen, $r$ sei ein Rest des Quadrats $aa$ nach dem Modul $m$, wenn $r\equiv aa\,(\operatorname{mod.}m)$ gilt. Der Kürze wegen nennen wir in diesem Abschnitt aber stets \emph{$r$} einen quadratischen Rest \emph{von $m$ selbst}, ohne dass dadurch eine Mehrdeutigkeit zu befürchten ist. Denn das Wort \emph{Rest} verwenden wir von nun an nicht mehr im Sinn einer kongruenten Zahl, außer wenn wir von \emph{kleinsten} Resten sprechen; dabei kann kein Zweifel entstehen.}, die letzteren ihre \emph{quadratischen Nichtreste}, oder auch einfach \emph{Reste} und \emph{Nichtreste}, wenn dadurch keine Mehrdeutigkeit entsteht. Es genügt, alle Zahlen $0,1,2\ldots m-1$ einzuordnen, denn kongruente Zahlen gehören zur selben Klasse.

Auch diese Untersuchung beginnen wir mit Primzahlmoduln; das ist auch dann zu verstehen, wenn wir es nicht ausdrücklich erwähnen. Die Primzahl $2$ schließen wir jedoch aus, betrachten also nur \emph{ungerade} Primzahlen.

\Needspace{10\baselineskip}
\begin{center}\small\emph{Ist der Modul eine Primzahl, so gibt es unter den kleineren Zahlen ebenso viele Reste wie Nichtreste.}\end{center}

\Article{96}
\emph{Wählt man die Primzahl $p$ als Modul, so ist die Hälfte der Zahlen $1,2,3\GaussFourDots p-1$ quadratischer Rest, die übrigen sind Nichtreste; es gibt also $\tfrac12(p-1)$ Reste und ebenso viele Nichtreste.}

Man zeigt leicht, dass alle Quadrate $1,4,9\ldots\tfrac14(p-1)^2$ inkongruent sind. Wäre nämlich $rr\equiv r'r'(\operatorname{mod.}p)$ möglich und wären die Zahlen $r,r'$ verschieden und nicht größer als $\tfrac12(p-1)$, so wäre, wenn wir $r>r'$ annehmen, was erlaubt ist, $(r-r')(r+r')$ positiv und \SourceLeaf{80}durch $p$ teilbar. Aber beide Faktoren $r-r'$ und $r+r'$ sind kleiner als $p$, sodass diese Annahme unmöglich ist (Artikel 13). Es gibt also $\tfrac12(p-1)$ quadratische Reste unter den Zahlen $1,2,3\GaussFourDots p-1$. Mehr kann es nicht geben, denn mit dem Rest $0$ erhält man $\tfrac12(p+1)$, und die Anzahl aller Reste kann diesen Wert nicht übersteigen. Die übrigen Zahlen sind daher Nichtreste, und ihre Anzahl ist $=\tfrac12(p-1)$.

Da null stets ein Rest ist, schließen wir diese Zahl und die durch den Modul teilbaren Zahlen aus unseren Untersuchungen aus. Dieser Fall ist für sich klar und würde nur die übersichtliche Form der Sätze stören. Aus demselben Grund haben wir auch den Modul $2$ ausgeschlossen.

\Article{97}
Mehrere Ergebnisse dieses Abschnitts lassen sich auch aus den Grundsätzen des vorigen Abschnitts herleiten. Da es nützlich ist, dieselbe Wahrheit mit verschiedenen Methoden zu untersuchen, wollen wir diesen Zusammenhang zeigen. Man sieht leicht, dass alle einem Quadrat kongruenten Zahlen \emph{gerade} Indizes haben, dagegen diejenigen, die keinem Quadrat kongruent sein können, ungerade. Da $p-1$ gerade ist, gibt es ebenso viele gerade wie ungerade Indizes, nämlich $\tfrac12(p-1)$, und damit ebenso viele Reste wie Nichtreste.

\emph{Beispiele.}\quad Für die folgenden Moduln sind die Reste
\begin{center}
\begin{tabular}{@{}r@{\quad}c@{\quad}l@{}}
$3$ & $\mathinner{\ldotp\ldotp\ldotp\ldotp\ldotp}$ & $1.$\\
$5$ & $\mathinner{\ldotp\ldotp\ldotp\ldotp\ldotp}$ & $1,4.$\\
$7$ & $\mathinner{\ldotp\ldotp\ldotp\ldotp\ldotp}$ & $1,2,4.$\\
$11$ & $\mathinner{\ldotp\ldotp\ldotp\ldotp\ldotp}$ & $1,3,4,5,9.$\\
$13$ & $\mathinner{\ldotp\ldotp\ldotp\ldotp\ldotp}$ & $1.3.4,9,10,12.$\\
$17$ & $\GaussFourDots,$ & $1,2,4,8,9,13,15,16.$\\
\multicolumn{3}{c}{usw.}
\end{tabular}
\end{center}
Die übrigen Zahlen unterhalb dieser Moduln sind Nichtreste.

\Needspace{10\baselineskip}
\begin{center}\small\emph{Ob eine zusammengesetzte Zahl Rest oder Nichtrest einer gegebenen Primzahl ist, hängt von ihren Faktoren ab.}\end{center}

\Article{98}
\textsc{Satz.}\quad \emph{Das Produkt zweier quadratischer Reste der Primzahl $p$ ist ein Rest; das Produkt eines Restes und eines Nichtrestes ist ein Nichtrest; das Produkt zweier Nichtreste ist ein Rest.}

\SourceLeaf{81}\emph{Beweis.}\quad I.\quad Seien $A,B$ Reste, die aus den Quadraten $aa,bb$ entstehen, also $A\equiv aa$, $B\equiv bb$. Dann ist das Produkt $AB$ dem Quadrat der Zahl $ab$ kongruent, also ein Rest.

II.\quad Ist $A$ ein Rest, etwa $\equiv aa$, und $B$ ein Nichtrest, so ist $AB$ ein Nichtrest. Nehmen wir nämlich, wenn möglich, $AB\equiv kk$ an und setzen $\frac{k}{a}\,(\mathrm{mod.}\,p)\equiv b$. Dann gilt $aaB\equiv aabb$, also $B\equiv bb$; somit wäre $B$ ein Rest, entgegen der Voraussetzung.

\emph{Anderer Beweis.}\quad Wir multiplizieren alle Reste unter den Zahlen $1,2,3\GaussFourDots p-1$ (ihre Anzahl ist $=\tfrac12(p-1)$) mit $A$. Alle Produkte sind quadratische Reste und untereinander inkongruent. Multiplizieren wir nun den Nichtrest $B$ mit $A$, so ist das Produkt keinem der schon erhaltenen Produkte kongruent. Wäre es ein Rest, so hätten wir $\tfrac12(p+1)$ inkongruente Reste, noch ohne den Rest $0$, entgegen Artikel 96.

III.\quad Seien $A,B$ Nichtreste. Multiplizieren wir alle Reste unter den Zahlen $1,2,3\GaussFourDots p-1$ mit $A$, so erhalten wir $\tfrac12(p-1)$ untereinander inkongruente Nichtreste (II). Das Produkt $AB$ kann keinem davon kongruent sein. Wäre es also ein Nichtrest, so hätten wir $\tfrac12(p+1)$ untereinander inkongruente Nichtreste, entgegen Artikel 96. Also ist das Produkt usw. \mbox{\emph{Q. E. D.}}

Noch leichter lassen sich diese Sätze aus den Grundsätzen des vorigen Abschnitts herleiten. Die Indizes der Reste sind stets gerade, die der Nichtreste ungerade. Der Index des Produkts zweier Reste oder zweier Nichtreste ist daher gerade, das Produkt selbst also ein Rest. Dagegen ist der Index des Produkts eines Restes und eines Nichtrestes ungerade und das Produkt selbst ein Nichtrest.

Beide Beweismethoden lassen sich auch auf die folgenden Sätze anwenden: \emph{Der Wert des Ausdrucks $\frac{a}{b}\,(\mathit{mod.}\,p)$ ist ein Rest, wenn die Zahlen $a,b$ beide Reste oder beide Nichtreste sind; dagegen ist er ein Nichtrest, wenn von den Zahlen $a,b$ eine ein Rest und die andere ein Nichtrest ist.} Diese Sätze ergeben sich auch durch Umkehrung der vorigen.

\Article{99}
Allgemein ist ein Produkt beliebig vieler Faktoren sowohl dann ein Rest, wenn alle Faktoren Reste sind, als auch dann, wenn unter ihnen eine gerade Anzahl von Nichtresten vorkommt. Ist die Anzahl der Nichtreste unter den Faktoren ungerade, so ist das Produkt ein Nichtrest. Man kann daher leicht entscheiden, \SourceLeaf{82}ob eine zusammengesetzte Zahl ein Rest ist, wenn man weiß, wie es sich mit jedem ihrer Faktoren verhält. Deshalb haben wir in Tabelle II nur Primzahlen aufgenommen. Die Tabelle ist folgendermaßen angeordnet: Am Rand stehen die Moduln\HistoricalNote{*)}{Wie wir auch auf zusammengesetzte Moduln verzichten können, werden wir bald zeigen.}, im Inneren die aufeinanderfolgenden Primzahlen. Ist eine dieser Primzahlen ein Rest eines Moduls, so steht im zugehörigen Feld ein kurzer Strich; ist sie ein Nichtrest des Moduls, so bleibt dieses Feld leer.

\Needspace{10\baselineskip}
\begin{center}\small\emph{Zusammengesetzte Moduln.}\end{center}

\Article{100}
Bevor wir zu schwierigeren Fragen übergehen, müssen wir noch einiges über nichtprime Moduln ergänzen.

Wählt man für eine Primzahl $p$ eine Potenz $p^n$ als Modul (wobei wir annehmen, dass $p$ nicht $2$ ist), so ist von allen Zahlen, die nicht durch $p$ teilbar und kleiner als der Modul sind, die eine Hälfte Rest, die andere Nichtrest. Die Anzahl beider ist also $=\tfrac12(p-1)p^{n-1}$.

Ist nämlich $r$ ein Rest, so ist er einem Quadrat kongruent, dessen Wurzel die Hälfte des Moduls nicht übersteigt; siehe Artikel 94. Man sieht leicht, dass es $\tfrac12(p-1)p^{n-1}$ Zahlen gibt, die nicht durch $p$ teilbar und kleiner als die Hälfte des Moduls sind. Es bleibt also zu zeigen, dass ihre Quadrate alle inkongruent sind, also verschiedene quadratische Reste liefern. Wären die Quadrate zweier Zahlen $a,b$, die nicht durch $p$ teilbar und kleiner als die Hälfte des Moduls sind, kongruent, so wäre $aa-bb$ oder $(a-b)(a+b)$ durch $p^n$ teilbar (wobei wir $a>b$ annehmen dürfen). Das ist nur möglich, wenn \emph{entweder} eine der Zahlen $a-b.\ a+b$ durch $p^n$ teilbar ist, was nicht sein kann, da beide $<p^n$ sind, \emph{oder} die eine durch $p^m$, die andere durch $p^{n-m}$, also beide durch $p$ teilbar sind. Aber auch das ist unmöglich. Denn dann wären ihre Summe und Differenz $2a$ und $2b$ durch $p$ teilbar, also auch $a$ und $b$, entgegen der Voraussetzung. --- Daraus folgt, dass es unter den nicht durch $p$ teilbaren Zahlen unterhalb des Moduls $\tfrac12(p-1)p^{n-1}$ Reste gibt; die übrigen, deren Anzahl ebenso groß ist, sind Nichtreste. \mbox{\emph{Q. E. D.}} --- Diesen Satz kann man auch wie in Artikel 97 durch Betrachtung der Indizes herleiten.

\Needspace{8\baselineskip}
\Article{101}
\emph{Jede nicht durch $p$ teilbare Zahl, die ein Rest von $p$ ist, ist auch ein Rest \SourceLeaf{83}von $p^n$; jeder Nichtrest von $p$ ist auch ein Nichtrest von $p^n$.}

Der zweite Teil dieser Aussage ist für sich klar. Wäre der erste falsch, so gäbe es unter den Zahlen, die kleiner als $p^n$ und nicht durch $p$ teilbar sind, mehr Reste von $p$ als von $p^n$, also mehr als $\frac{1}{2}p^{n-1}(p-1)$. Man sieht aber ohne Schwierigkeit, dass die Anzahl der Reste von $p$ unter diesen Zahlen genau $=\frac{1}{2}p^{n-1}(p-1)$ ist.

Ebenso leicht ist es, ein Quadrat zu finden, das einem gegebenen Rest nach dem Modul $p^n$ kongruent ist, wenn ein ihm nach dem Modul $p$ kongruentes Quadrat bekannt ist.

Hat man nämlich ein Quadrat $aa$, das dem gegebenen Rest $A$ nach dem Modul $p^\mu$ kongruent ist, so erhält man daraus ein Quadrat, das $A$ nach dem Modul $p^\nu$ kongruent ist (wobei wir $\nu>\mu$ und $=$ oder $<2\mu$ voraussetzen), folgendermaßen. Wir setzen die Wurzel des gesuchten Quadrats $=\pm a+xp^\mu$; diese Form muss sie offensichtlich haben. Dann muss $aa\pm2axp^\mu+xxp^{2\mu}\equiv A\,(\mathrm{mod.}\,p^\nu)$ gelten, oder wegen $2\mu\geq\nu$ auch $A-aa\equiv\pm2axp^\mu\,(\mathrm{mod.}\,p^\nu)$. Sei $A-aa=p^{\mu}d$; dann ist $x$ der Wert des Ausdrucks $\pm\frac{d}{2a}\,(\mathrm{mod.}\,p^{\nu-\mu})$, der dem Ausdruck $\pm\frac{A-aa}{2ap^\mu}\,(\mathrm{mod.}\,p^\nu)$ gleichwertig ist.

Ein Quadrat, das $A$ nach dem Modul $p$ kongruent ist, liefert also ein Quadrat, das $A$ nach dem Modul $p^2$ kongruent ist. Von hier kann man zum Modul $p^4$, dann zu $p^8$ usw. aufsteigen.

\emph{Beispiel.} Ist der Rest $6$ gegeben, der nach dem Modul $5$ dem Quadrat $1$ kongruent ist, so findet man das Quadrat $9^2$, dem er nach $25$ kongruent ist, ferner $16^2$, dem er nach $125$ kongruent ist, usw.

\Article{102}
Was die durch $p$ teilbaren Zahlen betrifft, so sind ihre Quadrate offensichtlich durch $pp$ teilbar. Daher sind alle Zahlen, die zwar durch $p$, aber nicht durch $pp$ teilbar sind, Nichtreste von $p^n$. Ist allgemein eine Zahl $p^k A$ gegeben, wobei $A$ nicht durch $p$ teilbar ist, so sind folgende Fälle zu unterscheiden:

1) Wenn $k=$ oder $>n$, so gilt $p^k A\equiv0\,(\mathrm{mod.}\,p^n)$, also handelt es sich um einen Rest.

2) Wenn $k<n$ und der Exponent ungerade ist, so ist $p^k A$ ein Nichtrest.

Wäre nämlich $p^k A=p^{2\kappa+1}A\equiv ss\,(\mathrm{mod.}\,p^n)$, so wäre $ss$ durch $p^{2\kappa+1}$ teilbar. Das ist nur möglich, wenn $s$ durch $p^{\kappa+1}$ teilbar ist. Dann wäre aber $ss$ \SourceLeaf{84}auch durch $p^{2\kappa+2}$ teilbar und somit (da $2\kappa+2$ sicher nicht größer als $n$ ist) auch $p^k A$, \emph{das heißt} $p^{2\kappa+1}A$. Also wäre $A$ durch $p$ teilbar, entgegen der Voraussetzung.

3) Wenn $k<n$ und der Exponent gerade ist. Dann ist $p^k A$ ein Rest oder Nichtrest von $p^n$, je nachdem, ob $A$ ein Rest oder Nichtrest von $p$ ist. Ist nämlich $A$ ein Rest von $p$, so ist es auch ein Rest von $p^{n-k}$. Setzen wir $A\equiv aa (\mathrm{mod.}\,p^{n-k})$, so gilt $Ap^k\equiv aap^k\,(\mathrm{mod.}\,p^n)$, und $aap^k$ ist ein Quadrat. Ist dagegen $A$ ein Nichtrest von $p$, so kann $p^k A$ kein Rest von $p^n$ sein. Denn nehmen wir $p^k A\equiv aa\,(\mathrm{mod.}\,p^n)$ an, so muss $aa$ durch $p^k$ teilbar sein. Der Quotient ist ein Quadrat, dem $A$ nach dem Modul $p^{n-k}$ und damit auch nach $p$ kongruent ist. \emph{Das heißt,} $A$ wäre ein Rest von $p$, entgegen der Voraussetzung.

\Article{103}
Da wir den Fall $p=2$ ausgeschlossen haben, müssen wir darüber noch einiges sagen. Ist $2$ der Modul, so ist jede Zahl ein Rest, und es gibt keine Nichtreste. Ist aber $4$ der Modul, so sind alle ungeraden Zahlen der Form $4k+1$ Reste, alle Zahlen der Form $4k+3$ Nichtreste. Ist schließlich $8$ oder eine höhere Potenz von $2$ der Modul, so sind alle ungeraden Zahlen der Form $8k+1$ Reste; die übrigen, also die Zahlen der Formen $8k+3,\ 8k+5,\ 8k+7$, sind Nichtreste. Der zweite Teil dieser Aussage ist klar, denn das Quadrat jeder ungeraden Zahl, ob sie die Form $4k+1$ oder $4k-1$ hat, hat die Form $8k+1$. Den ersten Teil beweisen wir so.

1) Ist die Summe oder die Differenz zweier Zahlen durch $2^{n-1}$ teilbar, so sind ihre Quadrate nach dem Modul $2^n$ kongruent. Setzt man nämlich die eine $=a$, so hat die andere die Form $2^{n-1}h\pm a$, deren Quadrat $\equiv aa\,(\mathrm{mod.}\,2^n)$ ist.

2) Jede ungerade Zahl, die ein quadratischer Rest von $2^n$ ist, ist einem Quadrat kongruent, dessen Wurzel ungerade und $<2^{n-2}$ ist. Sei $aa$ ein beliebiges Quadrat, dem jene Zahl kongruent ist, und sei $a\equiv\pm\alpha (\mathrm{mod.}\,2^{n-1})$, wobei $\alpha$ die Hälfte des Moduls nicht übersteigt (Artikel 4). Dann ist $aa\equiv\alpha\alpha$, also ist auch die gegebene Zahl $\equiv\alpha\alpha$. Offensichtlich sind sowohl $a$ als auch $\alpha$ ungerade, und es gilt $\alpha<2^{n-2}$.

\Needspace{5\baselineskip}
3) Die Quadrate aller ungeraden Zahlen unterhalb von $2^{n-2}$ sind nach $2^n$ inkongruent. Seien nämlich $r$ und $s$ zwei solche Zahlen. Wären ihre Quadrate nach $2^n$ kongruent, so wäre $(r-s)(r+s)$ durch $2^n$ teilbar (wobei $r>s$ angenommen ist). Man sieht aber leicht, dass die Zahlen $r-s,\ r+s$ nicht beide durch $4$ teilbar sein können. \SourceLeaf{85}Ist also die eine nur durch $2$ teilbar, so müsste, damit das Produkt durch $2^n$ teilbar wäre, die andere durch $2^{n-1}$ teilbar sein. \mbox{\emph{Q. E. A.}} Denn beide sind $<2^{n-2}$.

4) Reduzieren wir schließlich diese Quadrate auf ihre \emph{kleinsten positiven Reste}, so erhalten wir $2^{n-3}$ verschiedene quadratische Reste unterhalb des Moduls\HistoricalNote{*)}{Denn die Anzahl der ungeraden Zahlen unterhalb von $2^{n-2}$ ist $2^{n-3}$.}, von denen jeder die Form $8k+1$ hat. Da es aber genau $2^{n-3}$ Zahlen der Form $8k+1$ unterhalb des Moduls gibt, müssen sie alle unter diesen Resten vorkommen. \mbox{\emph{Q. E. D.}}

Um ein Quadrat zu finden, das einer gegebenen Zahl der Form $8k+1$ nach dem Modul $2^n$ kongruent ist, kann man ähnlich wie in Artikel 101 vorgehen; siehe auch Artikel 88. --- Für gerade Zahlen gilt schließlich dasselbe, was wir in Artikel 102 allgemein dargelegt haben.

\Article{104}
Über die Anzahl der verschiedenen Werte (\emph{das heißt,} der nach dem Modul inkongruenten Werte), die ein Ausdruck wie $V=\surd A\,(\mathrm{mod.}\,p^n)$ zulässt, wenn $A$ ein Rest von $p^n$ ist, ergibt sich aus dem Vorigen leicht Folgendes. (Wir nehmen wie zuvor an, dass $p$ eine Primzahl ist, und schließen der Kürze wegen den Fall $n=1$ gleich mit ein.) I. Ist $A$ nicht durch $p$ teilbar, so hat $V$ für $p=2,\ n=1$ \emph{einen} Wert, nämlich $V \equiv1$; \emph{zwei}, wenn $p$ ungerade ist, und ebenso für $p=2,\ n=2$: Setzt man einen $\equiv v$, so ist der andere $\equiv-v$; \emph{vier} für $p=2,\ n>2$: Setzt man einen $\equiv v$, so sind die übrigen $\equiv-v,\ 2^{n-1}+v,\ 2^{n-1}-v$. II. Ist $A$ durch $p$, aber nicht durch $p^n$ teilbar, so sei die höchste Potenz von $p$, die $A$ teilt, $p^{2\mu}$ (ihr Exponent muss offensichtlich gerade sein), und es sei $A=ap^{2\mu}$. Dann sind alle Werte von $V$ durch $p^\mu$ teilbar, und die entstehenden Quotienten sind Werte des Ausdrucks $V'=\surd a\,(\mathrm{mod.}\,p^{n-2\mu})$. Alle verschiedenen Werte von $V$ erhält man daher, indem man alle Werte des Ausdrucks $V'$ zwischen $0$ und $p^{n-\mu}$ mit $p^\mu$ multipliziert. Sie werden also dargestellt durch

\[
vp^\mu,\quad
vp^\mu+p^{n-\mu},\quad
vp^\mu+2p^{n-\mu}\ldots vp^\mu+(p^\mu-1)p^{n-\mu}
\]

wobei $v$ nacheinander alle \emph{verschiedenen} Werte des Ausdrucks $V'$ bezeichnet. Ihre Anzahl ist somit $p^\mu,\ 2p^\mu$ oder $4p^\mu$, je nachdem, ob die Anzahl der letzteren (nach Fall I) $1,\ 2$ oder $4$ ist. III. Ist $A$ durch $p^n$ teilbar, so setze man $n=2m$ oder $=2m-1$, je nachdem, ob der Exponent gerade oder ungerade ist. Man sieht leicht, dass alle durch $p^m$ teilbaren Zahlen und keine anderen Werte von $V$ sind. Alle verschiedenen Werte sind daher $0,p^m,2p^m\ldots(p^{n-m}-1)p^m$, und ihre Anzahl ist $p^{n-m}$.

\SourceLeaf{86}\Article{105}
Es bleibt der Fall, in dem der Modul $m$ aus mehreren Primzahlen zusammengesetzt ist. Sei $m=abc\ldots$, wobei $a,b,c$ usw. verschiedene Primzahlen oder Potenzen verschiedener Primzahlen bezeichnen. Ist $n$ ein Rest von $m$, so ist $n$ offensichtlich auch ein Rest von jedem der Moduln $a,b,c$ usw.; also ist $n$ sicher ein Nichtrest von $m$, wenn es ein Nichtrest einer der Zahlen $a,b,c$ usw. ist. Umgekehrt ist $n$, wenn es ein Rest von jeder der Zahlen $a,b,c$ usw. ist, auch ein Rest ihres Produkts $m$. Nehmen wir nämlich $n\equiv A^2$, $B^2,C^2$ usw. nach den Moduln $a,b,c$ usw. an. Bestimmen wir eine Zahl $N$, die den Zahlen $A,B,C$ usw. nach den Moduln $a,b,c$ usw. jeweils kongruent ist (Artikel 32), so gilt $n\equiv NN$ nach allen diesen Moduln und damit auch nach ihrem Produkt $m$. --- Man sieht leicht, dass auf diese Weise die Verbindung \emph{jedes} Wertes von $A$, also des Ausdrucks $\surd n\,(\mathrm{mod.}\,a)$, mit \emph{jedem} Wert von $B$, mit \emph{jedem} Wert von $C$ usw. einen Wert von $N$, also des Ausdrucks $\surd n\,(\mathrm{mod.}\,m)$, liefert; dass verschiedene Verbindungen verschiedene Werte von $N$ ergeben; und dass sämtliche Verbindungen sämtliche Werte liefern. Die Anzahl aller verschiedenen Werte von $N$ ist daher gleich dem Produkt der Anzahlen der Werte von $A,B,C$ usw., deren Bestimmung wir im vorigen Artikel gezeigt haben. --- Ist ferner ein Wert des Ausdrucks $\surd n\,(\mathrm{mod.}\,m)$, also von $N$, bekannt, so ist dieser auch ein Wert von allen Zahlen $A,B,C$ usw. Da man daraus nach dem vorigen Artikel alle übrigen Werte dieser Größen herleiten kann, lassen sich aus einem Wert von $N$ leicht alle übrigen bestimmen.

\emph{Beispiel.} Sei der Modul $315$ gegeben und gefragt, ob $46$ ein Rest oder Nichtrest davon ist. Die Primteiler von $315$ sind $3,5,7$, und $46$ ist ein Rest von jedem von ihnen, also auch ein Rest von $315$. Da ferner $46\equiv1$ und $\equiv64 (\mathrm{mod.}\,9)$; $\equiv1$ und $\equiv16\,(\mathrm{mod.}\,5)$; $\equiv4$ und $\equiv25\,(\mathrm{mod.}\,7)$ gelten, findet man die Wurzeln der Quadrate, denen $46$ nach dem Modul $315$ kongruent ist: $19,26,44,89,226$, $271,289,296$.

\Needspace{10\baselineskip}
\begin{center}
\small\itshape
Allgemeines Kriterium dafür, ob eine gegebene Zahl Rest oder Nichtrest einer gegebenen Primzahl ist.
\end{center}

\Article{106}
Aus dem Vorigen folgt: Wenn wir stets entscheiden können, ob eine gegebene \emph{Primzahl} Rest oder Nichtrest einer \emph{gegebenen Primzahl} ist, lassen sich alle übrigen Fälle darauf zurückführen. Für diesen Fall müssen wir also mit aller Sorgfalt sichere Kriterien suchen. Bevor wir diese Untersuchung beginnen, geben wir ein Kriterium aus dem vorigen Abschnitt an, das zwar in der Praxis \SourceLeaf{87}fast keinen Nutzen hat, wegen seiner Einfachheit und Allgemeinheit aber eine Erwähnung verdient.

\emph{Jede Zahl $A$, die nicht durch die Primzahl $2m+1$ teilbar ist, ist ein Rest oder Nichtrest dieser Primzahl, je nachdem, ob $A^m\equiv+1$ oder $\equiv-1\,(\operatorname{mod.}2m+1)$ gilt.}

Für den Modul $2m+1$ sei der Index der Zahl $A$ in einem beliebigen System $a$. Dann ist $a$ gerade, wenn $A$ ein Rest von $2m+1$ ist, dagegen ungerade, wenn $A$ ein Nichtrest ist. Der Index der Zahl $A^m$ ist aber $ma$, \emph{das heißt} $\equiv0$ oder $\equiv m\,(\operatorname{mod.}2m)$, je nachdem, ob $a$ gerade oder ungerade ist. Also ist $A^m$ im ersten Fall $\equiv+1$, im zweiten $\equiv-1\,(\operatorname{mod.}2m+1)$. Siehe Artikel 57 und 62.

{\rightskip=0pt plus 3em
\emph{Beispiel.}\quad $3$ ist ein Rest von $13$, weil $3^6\equiv1\,(\operatorname{mod.}13)$ gilt; dagegen ist $2$ ein Nichtrest von $13$, weil $2^6\equiv-1\,(\operatorname{mod.}13)$ gilt.\par}

Sobald die zu prüfenden Zahlen aber auch nur mäßig groß sind, ist dieses Kriterium wegen des gewaltigen Rechenaufwands völlig unbrauchbar.

\Needspace{10\baselineskip}
\begin{center}\emph{Untersuchungen über die Primzahlen, für die gegebene Zahlen Reste oder Nichtreste sind.}\end{center}

\Article{107}
Für einen gegebenen Modul lassen sich alle Zahlen, die seine Reste oder Nichtreste sind, sehr leicht angeben. Setzt man jene Zahl $=m$, so muss man die Quadrate bestimmen, deren Wurzeln die Hälfte von $m$ nicht übersteigen, oder Zahlen, die diesen Quadraten nach $m$ kongruent sind (für die Praxis gibt es noch bequemere Methoden). Dann sind alle Zahlen, die einer von diesen nach $m$ kongruent sind, Reste von $m$; alle Zahlen, die keiner davon kongruent sind, sind Nichtreste. --- Die umgekehrte Frage, \emph{zu einer gegebenen Zahl alle Zahlen anzugeben, von denen sie Rest oder Nichtrest ist}, erfordert jedoch eine viel tiefere Untersuchung. Dieses Problem, von dessen Lösung das im vorigen Artikel gestellte abhängt, werden wir nun untersuchen und dabei mit den einfachsten Fällen beginnen.

\Needspace{10\baselineskip}
\begin{center}\emph{Der Rest $-1$.}\end{center}

\Article{108}
\textsc{Satz.}\quad \emph{Für alle Primzahlen der Form $4n+1$ ist $-1$ ein quadratischer Rest, für alle Primzahlen der Form $4n+3$ dagegen ein Nichtrest.}

\emph{Beispiel.}\quad $-1$ ist ein Rest der Zahlen $5,\allowbreak 13,\allowbreak 17,\allowbreak 29,\allowbreak 37,\allowbreak 41,\allowbreak 53,\allowbreak 61.73,\allowbreak 89,\allowbreak 97$ usw., der sich jeweils \SourceLeaf{88}aus den Quadraten der Zahlen $2,\allowbreak 5,\allowbreak 4,\allowbreak 12,\allowbreak 6,\allowbreak 9,\allowbreak 23,\allowbreak 11,\allowbreak 27,\allowbreak 34,\allowbreak 22$ usw. ergibt; dagegen ist es ein Nichtrest der Zahlen $3,\allowbreak 7,\allowbreak 11,\allowbreak 19,\allowbreak 23,\allowbreak 31,\allowbreak 43,\allowbreak 47,\allowbreak 59,\allowbreak 67,\allowbreak 71,\allowbreak 79,\allowbreak 83$ usw.

Diesen Satz haben wir schon in Artikel 64 erwähnt. Sein Beweis folgt leicht aus Artikel 106. Für eine Primzahl der Form $4n+1$ gilt nämlich $(-1)^{2n}\equiv1$, für eine Zahl der Form $4n+3$ dagegen $(-1)^{2n+1}\equiv-1$. Dieser Beweis stimmt mit dem dort gegebenen überein. Wegen der Eleganz und des Nutzens des Satzes wird ein weiterer Beweis auf anderem Weg aber nicht überflüssig sein.

\Article{109}
Die Gesamtheit aller Reste der Primzahl $p$, die kleiner als $p$ sind, mit Ausnahme des Restes $0$, bezeichnen wir mit dem Buchstaben $C$. Da ihre Anzahl stets $=\frac{p-1}{2}$ ist, ist sie gerade, wenn $p$ die Form $4n+1$ hat, und ungerade, wenn $p$ die Form $4n+3$ hat. Wie in Artikel 77, wo es um Zahlen im Allgemeinen ging, nennen wir solche Reste \emph{Partner}, deren Produkt $\equiv1$ (mod.~$p$) ist. Denn ist $r$ ein Rest, so ist offensichtlich auch $\frac{1}{r}$ (mod.~$p$) ein Rest. Da derselbe Rest unter den Resten $C$ nicht mehrere Partner haben kann, lassen sich alle Reste $C$ in Klassen einteilen, von denen jede zwei als Partner zusammengehörige Reste enthält. Gäbe es keinen Rest, der sein eigener Partner ist, \emph{das heißt,} enthielte jede Klasse zwei \emph{verschiedene} Reste, so wäre die Anzahl aller Reste doppelt so groß wie die Anzahl aller Klassen. Gibt es dagegen Reste, die ihre eigenen Partner sind, \emph{das heißt} Klassen, die nur einen Rest oder, wenn man lieber will, denselben Rest zweimal enthalten, und setzt man die Anzahl dieser Klassen $=a$, die der übrigen $=b$, so ist die Anzahl aller Reste $C$ dann $=a+2b$. Hat also $p$ die Form $4n+1$, so ist $a$ gerade; hat $p$ die Form $4n+3$, so ist $a$ ungerade. Zahlen kleiner als $p$ können aber nur dann ihre eigenen Partner sein, wenn sie $1$ oder $p-1$ sind (siehe Artikel 77), und die erste Zahl, $1$, kommt sicher unter den Resten vor. Daher muss im ersten Fall $p-1$ (oder, was hier dasselbe bedeutet, $-1$) ein Rest sein, im zweiten dagegen ein Nichtrest; sonst wäre im ersten Fall $a=1$, im zweiten $=2$, was unmöglich ist.

\Article{110}
Auch dieser Beweis stammt von dem bedeutenden Euler, der auch als Erster den vorigen Beweis gefunden hat. Siehe \emph{Opusc. Anal. Bd. I, S. 135.} --- Man erkennt leicht, dass er auf ähnlichen Grundsätzen beruht wie unser zweiter Beweis des Wilsonschen Satzes in Artikel 77. Wenn wir \SourceLeaf{89}aber diesen Satz voraussetzen, lässt sich der Beweis noch leichter führen. Unter den Zahlen $1,2,3\ldots p-1$ gibt es nämlich $\frac{p-1}{2}$ quadratische Reste von $p$ und ebenso viele Nichtreste. Die Anzahl der Nichtreste ist daher gerade, wenn $p$ die Form $4n+1$ hat, und ungerade, wenn $p$ die Form $4n+3$ hat. Somit ist das Produkt aller Zahlen $1,2,3\ldots p-1$ im ersten Fall ein Rest, im zweiten ein Nichtrest (Artikel 99). Dieses Produkt ist aber stets $\equiv-1$ (mod.~$p$); also ist auch $-1$ im ersten Fall ein Rest, im zweiten ein Nichtrest.

\Article{111}
Ist also $r$ ein Rest einer Primzahl der Form $4n+1$, so ist auch $-r$ ein Rest dieser Primzahl. Alle Nichtreste einer solchen Zahl bleiben dagegen auch bei umgekehrtem Vorzeichen Nichtreste\HistoricalNote{*)}{Wenn wir also über eine Zahl als Rest oder Nichtrest einer Zahl der Form $4n+1$ sprechen, können wir ihr Vorzeichen ganz außer Acht lassen oder ihr auch das doppelte Vorzeichen $\pm$ geben.}. Bei Primzahlen der Form $4n+3$ ist es umgekehrt: Beim Wechsel des Vorzeichens werden Reste zu Nichtresten und Nichtreste zu Resten. Siehe Artikel 98.

Aus dem Vorigen ergibt sich außerdem leicht die allgemeine Regel: $-1$ \emph{ist ein Rest aller Zahlen, die weder durch $4$ noch durch eine Primzahl der Form $4n+3$ teilbar sind, und ein Nichtrest aller übrigen Zahlen.} Siehe Artikel 103 und 105.
